% MIPRO 2027 paper draft — main file.
% MIPRO uses the IEEE conference template (IEEEtran, two columns, max. 6 pages).
% Check the 2027 call for papers for the exact page limit and deadline.
% Build: pdflatex main && bibtex main && pdflatex main && pdflatex main

\documentclass[conference]{IEEEtran}
\usepackage[utf8]{inputenc}
\usepackage[T1]{fontenc}
\usepackage{amsmath}
\usepackage{graphicx}
\usepackage{url}

\begin{document}

\title{Graph or Embeddings? Retrieval for Question Answering\\over a Scholarly Knowledge Graph}

\author{\IEEEauthorblockN{David Josipović}
\IEEEauthorblockA{% TODO: affiliation
University, Faculty\\
City, Croatia\\
davidjosipovic19@gmail.com}
% TODO: mentor as second author?
}

\maketitle

\begin{abstract}
Retrieval-augmented generation over scholarly knowledge graphs is usually
evaluated with an LLM judge. We build a graded relevance gold set for 50
questions over the Open Research Knowledge Graph (ORKG) and answer metrics that
need no LLM judge, and use them to compare two retrieval strategies inside one
pipeline: graph-native SPARQL retrieval and dense retrieval over text cards of
all 68,806 ORKG papers. With context construction, prompt and LLM held fixed,
both reach similar retrieval quality (nDCG@8 0.75 vs.\ 0.71 on the test split),
but they retrieve largely different papers, and a rank fusion of the two
reaches 0.79. Graph retrieval is better on method comparisons and claim
verification, where it supplies the deciding evidence; dense retrieval is
better on topical and title lookups. End-to-end answer quality is equal for two
LLMs, and the difference is smaller than the variance between repeated runs.
Across nine LLMs, an LLM judge correlates only moderately with the gold metrics
($\rho = 0.45$) and overstates the gap between cloud and local models.
% TODO: trim to the MIPRO abstract limit after the paper is complete.
\end{abstract}

\begin{IEEEkeywords}
retrieval-augmented generation, knowledge graphs, ORKG, dense retrieval, evaluation, LLM-as-a-judge
\end{IEEEkeywords}

\section{Introduction}
\label{sec:intro}
% TODO (0.6 p.): KG-RAG over scholarly graphs as an alternative to vector RAG;
% such systems are evaluated with LLM judges that do not see the retrieved
% content; RQ1 is graph retrieval better than dense retrieval over the same
% data, RQ2 where do they differ, RQ3 does an LLM judge measure the
% difference; contributions.

\section{Related Work}
\label{sec:related}
% TODO (0.5 p.): GraphRAG / KG-RAG; text-to-SPARQL over ORKG (SciQA);
% dense retrieval and hybrids (RRF); RAG evaluation: LLM-as-a-judge (G-Eval,
% RAGAS), claim-level faithfulness (FActScore), known judge biases.
% Collect and verify every reference.

% Section III — System
% Draft 2026-10-20. Code: backend/rag/*, backend/kg/queries.py, eval/vector_baseline.py.

\section{System}
\label{sec:system}

Fig.~\ref{fig:pipeline} shows the pipeline. Both systems compared in this paper
share every step except the choice of the papers given to the LLM.

\begin{figure}[t]
\centering
% Fig. 1 — pipeline. Needs \usepackage{tikz} and
% \usetikzlibrary{positioning,arrows.meta} in main.tex.
% Fits one IEEE column (about 8.8 cm).

\begin{tikzpicture}[
    font=\footnotesize,
    node distance=4mm,
    >={Stealth[length=1.6mm]},
    step/.style={draw, rounded corners=1.5pt, align=center,
                 inner sep=2.5pt, minimum height=6mm},
    wide/.style={step, text width=7.6cm},
    half/.style={step, text width=3.55cm, minimum height=3.1cm, fill=gray!15},
    io/.style={align=center, inner sep=1pt},
]
  % shared input
  \node[io] (q) {Question};
  \node[wide, below=3mm of q] (an)
    {\textbf{Question analysis:} type (rules, then LLM),\\
     entities (LLM), normalisation to ORKG forms};

  % the only step that differs between the two systems
  \node[half, anchor=north east] (gr) at ([xshift=-1.5mm, yshift=-6mm]an.south)
    {\textbf{Graph retrieval}\\[1pt]
     SPARQL strategies per type\\
     (entities, comparison, claim\\
     evidence, benchmark, origin,\\
     title) $\to$ role-aware\\
     ranking and filtering $\to$\\
     enrichment of the top 15};
  \node[half, anchor=north west] (de) at ([xshift=1.5mm, yshift=-6mm]an.south)
    {\textbf{Dense retrieval}\\[1pt]
     embed the question\\
     (\texttt{mxbai-embed-large})\\
     $\to$ cosine similarity\\
     with cards of all\\
     68,806 ORKG papers};

  % shared output
  \node[wide, anchor=north] (top) at ([yshift=-5mm]an.south |- gr.south)
    {\textbf{Top 8 papers}};
  \node[wide, below=of top] (ctx)
    {\textbf{Context:} full ORKG records fetched live,\\
     up to 15 \emph{predicate: value} rows per paper};
  \node[wide, below=of ctx] (llm)
    {\textbf{LLM} with a prompt per question type\\
     (local via Ollama, or a cloud API)};
  \node[io, below=3mm of llm] (ans)
    {Answer citing papers by exact title\\
     (claims: SUPPORTED / REFUTED / INSUFFICIENT)};

  % arrows
  \draw[->] (q) -- (an);
  \draw[->] (an.south -| gr.north) -- (gr.north);
  \draw[->, dashed] (an.south -| de.north) -- (de.north)
    node[midway, right, font=\scriptsize] {question text};
  \draw[->] (gr.south) -- (gr.south |- top.north);
  \draw[->] (de.south) -- (de.south |- top.north);
  \draw[->] (top) -- (ctx);
  \draw[->] (ctx) -- (llm);
  \draw[->] (llm) -- (ans);
\end{tikzpicture}

\caption{The question-answering pipeline. Only the retrieval step (shaded)
differs between the two systems; dense retrieval uses only the question text
(dashed).}
\label{fig:pipeline}
\end{figure}

\subsection{Question Analysis}

The question is first assigned one of six types (Section~\ref{sec:method}).
Unambiguous forms are decided by rules (e.g.\ ``find the paper titled~\ldots'',
``compare~\ldots{} versus~\ldots'', a yes/no question stating a claim); the rest
are classified by the LLM. In parallel, the LLM extracts typed entities:
methods, datasets, tasks, research fields and metrics. If none are found,
keyword phrases from the question are used instead. A normalisation step
expands each entity to the surface forms used in ORKG, e.g.\ \emph{CNN} to
\emph{convolutional neural network}, and singular and plural forms.

\subsection{Graph Retrieval}

Graph retrieval queries the public ORKG SPARQL endpoint and uses no embeddings.
The question type selects a set of query strategies, each following a
different path in the graph from a paper through its contributions:

\begin{itemize}
\item \textbf{Entity paths}: contributions whose method, dataset or research
  problem matches an entity, and their combination when both a method and a
  dataset are given;
\item \textbf{Comparison}: papers whose contributions contain both compared
  methods;
\item \textbf{Claim evidence}: contributions covering most terms of the
  question, returned with the predicate that links each value (e.g.\
  \emph{baseline} versus \emph{proposed model});
\item \textbf{Benchmark traversal}: from a contribution through its model and
  benchmark nodes to the evaluated dataset and the reported scores;
\item \textbf{Origin lookup}: the earliest papers recording a method, as
  candidates for the paper that introduced it;
\item \textbf{Title lookup}: exact and keyword title search for paper-lookup
  questions.
\end{itemize}

All planned queries run in parallel. Their results are reduced to one
profile per paper and scored by the roles in which the question's entities
occur. For example, a method in a method or baseline role scores higher than
the same string elsewhere, an exact title match dominates, and papers that
carry benchmark results get a bonus. Short terms must match whole words, so
\emph{GAN} does not match \emph{organ}. When a question names both a method and
a dataset, papers matching only one of them are filtered out. The 15
best-scored candidates are then \emph{enriched}: their full contribution
records are fetched in one query and the papers are re-scored on the complete
record. This step matters because the matching query returns only the matched
value, so a paper that compares ResNet with VGG may first arrive with its VGG
row only. The top 8 papers are passed on.

\subsection{Dense Retrieval Baseline}

The dense baseline replaces only the step above. Every paper in the ORKG RDF dump
(68,806 papers) is represented by one text card: its title, research field
and the labels of its contribution values, research problems first, at most
1,500 characters. The cards are embedded with \texttt{mxbai-embed-large};
the question is embedded with the model's retrieval instruction, and the 8
papers with the highest cosine similarity are returned. The cards use the
same graph data that graph retrieval can reach, so the comparison is between
two ways of selecting papers from the same information, not between different
corpora.

\subsection{Context and Generation}

For the 8 selected papers, from either retriever, the full ORKG records are
fetched live and rendered as one block per paper: the title followed by
\emph{predicate: value} rows, up to 15 rows per paper. The rows are chosen in order of priority: values matching the question's
entities, benchmark results, then values under key predicates (method,
baseline, dataset, metric, research problem). Each row keeps the predicate,
so the LLM sees the role of a value, e.g.\ \emph{Baseline: Faster R-CNN}. Claim questions get an evidence-oriented layout of the same
rows. Papers are identified by title rather than by number, which prevents
citations of non-existent sources. A prompt template per question type asks
the LLM to answer only from the context, to name papers by their exact title,
and, for claims, to give a verdict (SUPPORTED, REFUTED or INSUFFICIENT) with
evidence. The system runs with local models through Ollama or with cloud APIs,
behind one client interface.

% Section IV — Evaluation Methodology
% Draft 2026-10-20. Numbers: eval/ANSWER_EVAL_REPORT.md, eval/RETRIEVAL_REPORT.md.
% References: references.bib (verified 2026-10-20).

\section{Evaluation Methodology}
\label{sec:method}

\subsection{Questions}

We use 50 natural-language questions over ORKG in six types:
topic search (9), method usage (9), dataset search (8), method comparison (8),
claim verification (8) and paper lookup (8); 14 are easy, 20 medium and 16 hard.
Questions with odd IDs form the development split and those with even IDs
the test split (25 each). All retrieval development (strategies, ranking,
query-classification rules) was done on the development split only; we report
results on all 50 questions and on the test split.

\subsection{Relevance Gold Set}

For every question we wrote a grading criterion and graded candidate papers
on a three-level scale: 2 (the paper answers the question, e.g.\ it reports
both compared methods on the named task), 1 (partially relevant, e.g.\ one of
the two methods, or the right method on another task) and 0 (not relevant).
Candidates were pooled from every retrieval variant run during development,
from both systems compared in this paper (including a second embedding model),
and from all papers cited in generated answers. The final set contains 2,633
judgments (757 of grade 2, 890 of grade 1, 986 of grade 0); every paper in the
top 8 of every reported system is judged.

Grades were proposed by an LLM (Claude) from each paper's ORKG record and
reviewed by the author. To estimate their reliability, the author re-graded a
stratified sample of 120 question--paper pairs (40 per grade) without seeing
the proposed grades. Agreement was 0.85 (Cohen's $\kappa = 0.78$, quadratic-weighted
$\kappa = 0.85$; $\kappa = 0.91$ for the binary decision ``grade 2 or not''),
and $\kappa = 0.91$ on the pairs retrieved only by the dense baseline, so the
added grades do not favour either system. The author had seen some of the
original grades before, so this is not a fully independent annotation.

For the eight claim-verification questions we additionally fixed a gold
verdict (SUPPORTED, REFUTED or INSUFFICIENT) from the ORKG records, with the
set of verdicts accepted as correct; one question whose evidence is mixed is
excluded from verdict scoring.

\subsection{Retrieval Metrics}

Each system returns 8 papers. We report nDCG@8 with linear gains (2/1/0) and the
ideal ranking computed from all judged papers of the question, together with
precision, recall of grade-2 papers and hit rate at 8. Differences between
systems are tested with a paired bootstrap over questions (5,000 resamples,
95\,\% percentile intervals).

Graph retrieval depends on the LLM that classifies the question and extracts
entities. We therefore report it in two modes: \emph{oracle}, with the gold
question type and hand-written entities, which isolates the retrieval
strategies, and \emph{end-to-end}, with entities extracted by \texttt{qwen2.5:7b}.
Dense retrieval does not use these steps.

\subsection{Answer Metrics without an LLM Judge}

Answers follow a fixed template and name papers by their exact title. A context
paper counts as \emph{cited} when its title appears in the answer body; the
final list of sources is ignored, because models tend to copy all context
papers into it. Per answer we compute:

\begin{itemize}
\item \textbf{CitePrec}: mean gold grade of the cited papers divided by 2;
\item \textbf{CiteUtil}: share of the grade-2 papers in the context that the answer cites;
\item \textbf{False abstention}: the answer cites nothing (or returns INSUFFICIENT)
  although the context contains a grade-2 paper.
\end{itemize}

The \textbf{gold Score} of an answer is $F_1$(CitePrec, CiteUtil) when the
context contains a grade-2 paper; otherwise it is 1 for an abstention and
CitePrec (at most 0.5) for an answer. Claim questions are scored by the
verdict, where INSUFFICIENT is also accepted when
none of the deciding papers was retrieved. This score is appropriate for
comparing LLMs on the same retrieval, but it rewards abstention when retrieval
fails, so a weaker retriever can obtain a higher Score. To compare retrieval
systems we therefore use a \textbf{strict end-to-end score} (E2E): no credit for
abstention, claim questions scored only by the gold verdict, and $F_1$(CitePrec,
CiteUtil) otherwise.

\subsection{Claim-Level Faithfulness}

Following claim-decomposition approaches \cite{min2023factscore,es2024ragas}, a
local judge (\texttt{qwen3:8b} via Ollama, not one of the evaluated models)
splits each answer body into at most 25 atomic claims and checks every claim
separately against the full context the model saw, labelling it supported (S),
unsupported (U) or contradicted (C). For S and C the judge must quote the
deciding context line verbatim; if the quote is not found in the context the
label becomes U. Faithfulness is the share of S claims per answer, averaged per
model.

The judge was validated on 100 random claims labelled blind by a stronger
model, with all 21 disagreements reviewed by the author. Accuracy was 0.79
($\kappa = 0.51$) for S/U/C and 0.84 ($\kappa = 0.60$) for supported versus not
supported. Errors were balanced (7 too strict, 7 too lenient), so per-model
means are unbiased, but single labels are only moderately reliable. The judge
rarely uses C and labels misattributions as U, so we report U and C together as
\emph{unsupported}.

\subsection{LLM-Judge Baseline}

For comparison with common practice \cite{liu2023geval,zheng2023judging}, every
answer was also scored by an LLM judge on groundedness, relevance and
completeness (0--5). The judge receives the question, the answer, the titles of
the retrieved papers and answer hints, but not the context itself. It is the
same \texttt{qwen3:8b} model as the faithfulness judge, so differences between
the two reflect the evaluation protocol rather than the judge model.

\subsection{Experimental Setup}

The two retrieval systems share the whole pipeline except the choice of the 8
papers (Section~\ref{sec:system}): the same full ORKG records, context rows,
prompt and LLM. Dense retrieval embeds one text card per paper (title, research
field and contribution values, at most 1,500 characters) for all 68,806 papers
in the ORKG RDF dump of 18 October 2026 with \texttt{mxbai-embed-large}
\cite{mxbai2024}, and ranks by cosine similarity; to test sensitivity to the
embedder we repeat retrieval with \texttt{nomic-embed-text} \cite{nussbaum2024nomic}.
A reciprocal rank fusion (RRF, $k = 60$) \cite{cormack2009rrf} of the two ranked lists
is computed at the retrieval level only.

Nine LLMs answer with graph retrieval: five local models via Ollama
(\texttt{qwen2.5:7b}, \texttt{llama3}, \texttt{llama3.2}, \texttt{mistral},
\texttt{phi3:mini}; context window 8,192 tokens) and four cloud models via the
Groq free tier (\texttt{gpt-oss-20b}, \texttt{gpt-oss-120b}, \texttt{qwen3.8-27b},
\texttt{allam-2-7b}). The graph--dense comparison of answers uses one local and
one cloud model (\texttt{qwen2.5:7b} and \texttt{gpt-oss-20b}); \texttt{qwen2.5:7b}
was run twice with each retriever to measure run-to-run variance.

% Section V — Results
% Draft 2026-10-20. Sources: eval/results/retrieval_*.json, benchmark_*.json,
% faithfulness_*.json; summarised in eval/ANSWER_EVAL_REPORT.md.

\section{Results}
\label{sec:results}

\subsection{Retrieval: Graph versus Dense}

\begin{table}[t]
\caption{Retrieval quality, nDCG@8. Graph retrieval in oracle mode; best of graph and \texttt{mxbai} in bold.}
\label{tab:retrieval}
\centering
\setlength{\tabcolsep}{4pt}
\begin{tabular}{lrccc}
\hline
Question type & $n$ & Graph & Dense (mxbai) & Dense (nomic) \\
\hline
Dataset search       & 8 & \textbf{0.817} & 0.688 & 0.259 \\
Method comparison    & 8 & \textbf{0.660} & 0.535 & 0.166 \\
Claim verification   & 8 & \textbf{0.657} & 0.563 & 0.340 \\
Method usage         & 9 & 0.710 & \textbf{0.817} & 0.532 \\
Paper lookup         & 8 & 0.583 & \textbf{0.724} & 0.247 \\
Topic search         & 9 & 0.850 & \textbf{0.979} & 0.866 \\
\hline
All                  & 50 & 0.716 & 0.725 & 0.413 \\
Test split           & 25 & 0.754 & 0.712 & 0.435 \\
\hline
\multicolumn{5}{p{0.93\columnwidth}}{\footnotesize Test split: graph end-to-end (\texttt{qwen2.5:7b}) 0.695; RRF of graph and mxbai 0.792.}
\end{tabular}
\end{table}

Table~\ref{tab:retrieval} compares the two retrievers. Overall they are on par:
0.716 versus 0.725 nDCG@8 on all questions, with a paired bootstrap interval of
the difference of $[-0.10, 0.09]$. During development, graph retrieval improved
from 0.527 to 0.754 on the test split, measured on the same final gold set, so
the parity is with a tuned graph retriever.

The per-type results are not on par. Graph retrieval is better where the
question has structure that maps onto graph relations: datasets used by a
paper, papers reporting both compared methods, and benchmark results that
decide a claim. Dense retrieval is better on broad topics, on method usage and on finding a
paper by its title, where textual similarity is sufficient and graph retrieval
depends on how an entity is spelled in ORKG.

The two systems retrieve largely different papers: on average only 0.42 of the
8 papers are shared. A reciprocal rank fusion of the two lists reaches 0.792 on
the test split, above both. In end-to-end mode graph retrieval loses 0.06
(0.754 to 0.695) through misclassified questions and wrongly extracted entities,
a source of error that dense retrieval does not have. Finally, the dense result
depends strongly on the embedder: with \texttt{nomic-embed-text} nDCG@8 drops to
0.413, and to 0.166 on method comparisons. The \texttt{mxbai} baseline is thus a
strong one.

\subsection{Answers: Graph versus Dense}

\begin{table}[t]
\caption{Answers with the same LLM and only the retriever changed (all 50 questions). E2E: strict end-to-end score; Cl.: grade-2 papers cited per claim question; MC: E2E on method comparison. For \texttt{qwen2.5:7b}, run 1 / run 2. Faithful for \texttt{gpt-oss-20b} on the 49 answers judged with both retrievers.}
\label{tab:answers}
\centering
\setlength{\tabcolsep}{2.5pt}
\begin{tabular}{llccccc}
\hline
Model & Retr. & E2E & Score & Faithful & Cl. & MC \\
\hline
qwen2.5 & graph & .748/.696 & .758/.726 & .771/.758 & 2.25/2.38 & .64/.60 \\
        & dense & .731/.733 & .809/.811 & .739/.733 & 0.38/1.25 & .37/.51 \\
gpt-oss-20b & graph & .745 & .830 & .688 & 2.12 & .68 \\
            & dense & .745 & .841 & .640 & 1.38 & .42 \\
\hline
\end{tabular}
\end{table}

Table~\ref{tab:answers} shows why the lenient gold Score cannot compare
retrievers: it is higher with dense retrieval for both models (e.g.\ 0.809
versus 0.758 for \texttt{qwen2.5:7b}), although dense retrieval supplies fewer
deciding papers for claim questions. The difference comes from accepted
abstentions when the evidence is missing.

With the strict E2E score the overall answer quality is equal. The graph minus
dense difference is $-0.01$ $[-0.10, 0.08]$ for \texttt{qwen2.5:7b}, pooled over
both runs, and $0.00$ $[-0.11, 0.11]$ for \texttt{gpt-oss-20b}. The run-to-run
variance is of the same size as the difference between the systems. Between the
two \texttt{qwen2.5:7b} runs the retrieved papers were identical for 49 of 50
questions with graph retrieval and for all 50 with dense retrieval. Still, the
graph E2E moved from 0.748 to 0.696 and the sign of the difference flipped, and
a single answer moved by 0.08--0.09 on average. Differences of a few points from
a single run should therefore not be read as a difference between systems.

The per-type pattern from retrieval carries over to the answers. On method
comparisons graph retrieval leads in all three measurements (0.60--0.68 versus
0.37--0.51). On claim and comparison questions, answers based on graph
retrieval cite 1.5 to 6 times more grade-2 papers. Verdict accuracy on claim
questions is nevertheless not higher, because the models often decide claims
from prior knowledge rather than from the retrieved evidence. On paper lookup,
dense retrieval leads (0.80 versus 0.61 for \texttt{gpt-oss-20b}).

Answers based on graph retrieval are more faithful to the context in all three
measurements, by 0.02 to 0.05. None of these differences is significant on its
own ($+0.03$ $[-0.03, 0.09]$ for \texttt{qwen2.5:7b}, $+0.05$ $[-0.04, 0.13]$ for
\texttt{gpt-oss-20b}), but the direction is consistent.

\subsection{Nine LLMs on Graph Retrieval}

\begin{table}[t]
\caption{Nine LLMs with graph retrieval (50 questions). Score: gold Score; Test: test split; Uns.: share of unsupported claims; Judge: LLM judge, overall 0--5.}
\label{tab:models}
\centering
\setlength{\tabcolsep}{4pt}
\begin{tabular}{lccccc}
\hline
Model & Score & Test & Faithful & Uns. & Judge \\
\hline
gpt-oss-20b      & \textbf{0.830} & \textbf{0.833} & 0.694 & 0.278 & 4.57 \\
gpt-oss-120b     & 0.798 & 0.745 & 0.722 & 0.246 & 4.53 \\
qwen3.8-27b      & 0.790 & 0.782 & \textbf{0.814} & 0.183 & 4.31 \\
qwen2.5:7b$^\ast$ & 0.758 & 0.779 & 0.771 & 0.218 & 3.46 \\
mistral$^\ast$   & 0.680 & 0.706 & 0.704 & 0.292 & 3.41 \\
llama3$^\ast$    & 0.669 & 0.705 & 0.684 & 0.315 & 3.27 \\
phi3:mini$^\ast$ & 0.663 & 0.656 & 0.703 & 0.284 & 3.49 \\
llama3.2$^\ast$  & 0.648 & 0.676 & 0.668 & 0.332 & 3.39 \\
allam-2-7b       & 0.167 & 0.180 & 0.195 & 0.846 & 1.17 \\
\hline
\multicolumn{6}{l}{\footnotesize $^\ast$ local (Ollama); others via Groq.}
\end{tabular}
\end{table}

Table~\ref{tab:models} reports nine LLMs on the same graph retrieval. With good
retrieval, the best local model, \texttt{qwen2.5:7b}, is within 0.07 of the best
cloud model. On the test split it is above \texttt{gpt-oss-120b} (0.779 versus
0.745), and the larger \texttt{gpt-oss-120b} does not beat \texttt{gpt-oss-20b}.
Given the run-to-run variance measured above (up to 0.05), the order of adjacent
models in Table~\ref{tab:models} is not reliable; the separation between groups
is. The models differ in style: \texttt{qwen3.8-27b} cites nearly all relevant
context papers (CiteUtil 0.98) and is the most faithful, while
\texttt{gpt-oss-20b} selects papers more precisely (CitePrec 0.83) but adds
more content that is not in the context. \texttt{allam-2-7b} is not usable for
the task: it abstains on 78\,\% of answerable questions, and its small context
window causes request failures on some questions.

Contradictions of the context are rare (1--3\,\% of claims, an underestimate
given the judge's behaviour). The main faithfulness problem is addition:
18--33\,\% of claims are not supported by the context. These are mostly
background knowledge, method explanations and results attributed to the wrong
paper.

\subsection{LLM Judge versus Gold Metrics}

The LLM judge agrees with the gold Score only moderately at the answer level
(Spearman $\rho = 0.45$, $n = 450$). At the model level the agreement is high
($\rho = 0.87$), but only because the judge separates cloud from local models;
within the five local models it is $\rho = 0.25$. The judge also overstates the
gap between them: \texttt{gpt-oss-20b} and \texttt{qwen2.5:7b} differ by 0.07 in
gold Score but by 1.1 points (22\,\% of the scale) in the judge's score, and
\texttt{phi3:mini}, second to last among local models by gold Score, is ranked
first among them by the judge. Because the judge sees only source titles, it
cannot reward an answer for citing the right papers or penalise it for
misusing them; it appears to reward style and length instead.

Faithfulness, judged by the same model with the context, measures a
different property from the gold Score: the two correlate at $\rho = 0.67$
across models and only $\rho = 0.15$ across answers. Choosing the right papers
and staying within their content are thus separate abilities, and an
evaluation that reports only one of them misses the other.


\section{Discussion and Limitations}
\label{sec:discussion}
% TODO (0.6 p.): why graph wins comparisons (SPARQL finds papers with both
% methods and benchmark scores) and dense wins topics and lookups; run-to-run
% variance as a methodological point; limitations: 50 questions (25 test),
% LLM-proposed relevance grades with author review, repeated run only for
% qwen2.5, hybrid only at retrieval level, ORKG coverage, Groq free tier.

\section{Conclusion}
\label{sec:conclusion}
% TODO (0.2 p.)

% \section*{Acknowledgment}

\bibliographystyle{IEEEtran}
\bibliography{references}

\end{document}
